\textbf{Contrastive SSL.} Contrastive predictive coding~\cite{oord2018representation} introduced the InfoNCE objective that underlies most contrastive methods. SimCLR~\cite{chen2020simple} learns by maximising agreement between two augmented views with a normalised-temperature cross-entropy loss and a projection head; MoCo~\cite{he2019momentum} uses a momentum-encoder queue of negatives and BYOL~\cite{grill2020bootstrap} removes negatives altogether. We reimplement only the SimCLR-style loss, at a scale of one small encoder and one unlabeled pool.

\textbf{Pretext tasks.} Rotation prediction~\cite{gidaris2018unsupervised} asks a network to recognise which of four rotations was applied. Kolesnikov et al.~\cite{kolesnikov2019revisiting} show that the quality of pretext-task representations depends strongly on architecture choices, which motivates our use of an identical encoder and probe for all systems. Su et al.~\cite{su2019boosting} use self-supervised auxiliary tasks (rotation and jigsaw) for few-shot learning, a different use from our pretrain-then-probe setting.

\textbf{Label efficiency.} S4L~\cite{zhai2019s} and SimCLRv2~\cite{chen2020big} combine SSL with few labels at ImageNet scale. Konstantakos et al.~\cite{konstantakos2024self} compare SSL methods when pretraining data are scarce, the closest work to ours; we differ in using an even smaller dataset, a classical PCA baseline and a label-budget sweep. The optimiser for all networks is Adam~\cite{kingma2014adam}.
